\documentclass{article}
\usepackage[T1]{fontenc}
\usepackage{lmodern}
\usepackage{graphicx}
\usepackage{amsmath,amssymb}
\usepackage[a4paper,margin=2cm]{geometry}
\usepackage[absolute,overlay]{textpos}
\setlength{\TPHorizModule}{1mm}
\setlength{\TPVertModule}{1mm}
\usepackage[indonesian]{babel}
\usepackage{hyperref}
\setlength{\emergencystretch}{2em}
% unit: Halaman 57

\begin{document}

% === Halaman 57 (python/native) ===

Contoh 2:  Misalkan a = 01010111 = x6 + x4 + x2 + x + 1  dan b = 10000011 = x7 + x + 1 


   a dan b  GF(28) 

  (dalam notasi Hex, a = ’57’ dan b = ’83’)

(i)

a + b = ’57’ + ’83’  =  (x6 + x4 + x2 + x + 1 ) +  (x7 + x + 1 ) 




= x7 + x6 + x4 + x2 + 2x + 2 (mod 2)



             = x7 + x6 + x4 + x2 + 0x  + 0




= x7 + x6 + x4 + x2 = 11010100 = ‘D4’


Dalam representasi biner:


 01010111


 10000011 +


 11010100  → sama dengan hasil operasi XOR


 a + b = 11010100 = a XOR b 

Bagi  tiap koefisien dengan 2,

lalu ambil sisanya

\end{document}
